\section{Instalación de {\color{brown}\texttt{yay}}}
\subsection{Proceso de instalación}
\begin{frame}
	\begin{center}
		{\color{blue}Pasos a seguir:}\vspace{3mm}
		
      \begin{tcolorbox}[title=Descargamos los ficheros desde \textbf{Github}:]

			\orden{git clone https://aur.archlinux.org/yay.git}
      \end{tcolorbox}\pause
			
      \begin{tcolorbox}[title=Ahora vamos al directorio y compilamos:]
			\orden{cd yay}\vspace{1mm}
      
			\orden{makepkg -si}
      \end{tcolorbox}\pause
			
		\visible<4->{Para instalar software solo hay que escribir {\color{red}\texttt{yay}} en lugar de {\color{red}\texttt{pacman}} sin necesidad de usar la orden {\color{red}\texttt{sudo}}}.\\\visible<5>{Desde este momento podemos instalar todos los paquetes con {\color{red}\texttt{yay}}.}
		\end{center}
\end{frame}
